\documentclass[11pt]{article}

\usepackage[margin=1in]{geometry}
\usepackage{amsmath,amssymb,amsthm,mathtools}
\usepackage{graphicx}
\usepackage[hidelinks]{hyperref}

\newtheorem{theorem}{Theorem}
\newtheorem{proposition}{Proposition}
\newtheorem{corollary}{Corollary}
\newtheorem{remark}{Remark}

\newcommand{\T}{\mathbb T}
\newcommand{\Rd}{C_{\mathrm{rd}}}
\newcommand{\Deltas}{\,\Delta s}
\newcommand{\muT}{\mu}

\title{Non-Product Kamke $\Delta$-Functions and Redundancy of the Initial Derivative Clause}
\author{Automated Math Discovery Candidate Proof}
\date{August 11, 2026}

\begin{document}
\maketitle

\begin{center}
\fbox{\parbox{0.92\textwidth}{\textbf{Status.}
Candidate full answer to the Kamke $\Delta$-function question in the source
paper, likely valid pending expert review.  A broad polynomial family supplies
non-product examples.  More strongly, an example for which the initial
$\Delta$-derivative requirement is essential cannot exist under the stated
axioms.}}
\end{center}

\section*{Source and Question}

This packet concerns the final paragraph on Kamke $\Delta$-functions in
Section~5, page~29, of Oberta~\cite{Oberta}, arXiv:2512.13602v2.  Definition~7
there calls
\[
 w:[a,b]_{\T}\times[0,\infty)\longrightarrow[0,\infty)
\]
a Kamke $\Delta$-function when:
\begin{enumerate}
\item $w(\cdot,x)$ is rd-continuous for every $x\geq0$;
\item the family $\{w(t,\cdot):t\in[a,b]_{\T}\}$ is uniformly
equicontinuous on every bounded interval $[0,R]$;
\item $w(t,0)=0$;
\item zero is the unique nonnegative continuous solution of
\begin{equation}\label{eq:ineq}
 u(t)\leq \int_a^t w(s,u(s))\Deltas
\end{equation}
among those $u$ that are $\Delta$-differentiable at $a$ and satisfy
$u^\Delta(a)=0$.
\end{enumerate}
The paper asks for more sophisticated examples, for instance one not of the
form $q(t)h(x)$, or one having a nonzero continuous solution of
\eqref{eq:ineq} that is excluded only by the derivative requirement.

\begin{center}
\includegraphics[width=\textwidth]{figures/open_problem_crop.png}
\end{center}

\section*{Proof Intuition}

There are two independent observations.  First, the derivative condition in
Definition~7 is already forced by the integral inequality.  The inequality at
$t=a$ gives $u(a)=0$.  If $a$ is right-scattered, the first $\Delta$-integral
step is $\muT(a)w(a,u(a))=0$, so $u(\sigma(a))=0$.  If $a$ is right-dense,
uniform equicontinuity of all the functions $w(t,\cdot)$ at zero makes the
integrand uniformly small on the shrinking interval $[a,t]_{\T}$.  Dividing
by $t-a$ then shows $u(t)/(t-a)\to0$.  Thus the extra derivative condition
can never separate a nonzero solution from zero.

Second, many non-product examples are nevertheless easy to build.  Any finite
polynomial in $x$ with bounded nonnegative rd-continuous coefficients is at
most linear in $u$ on the bounded range of a candidate solution.  Time-scale
Gronwall therefore forces that solution to vanish.  Choosing nonproportional
coefficient functions prevents a factorization $q(t)h(x)$.

\section*{Automatic Endpoint Derivative}

We use the standard notation $\sigma$ for the forward jump and
$\muT(t)=\sigma(t)-t$ for graininess.  As in the source paper, the compact
interval $[a,b]_{\T}$ is treated as the ambient time scale and $a<b$.

\begin{theorem}[The derivative clause is automatic]\label{thm:auto}
Suppose $w$ satisfies axioms (i)--(iii) of Definition~7 in the source paper.
If a nonnegative continuous function $u:[a,b]_{\T}\to[0,\infty)$ satisfies
\eqref{eq:ineq} for every $t\in[a,b]_{\T}$, then $u$ is
$\Delta$-differentiable at $a$ and
\[
                              u^\Delta(a)=0.
\]
\end{theorem}

\begin{proof}
Putting $t=a$ in \eqref{eq:ineq} gives $0\leq u(a)\leq0$, hence $u(a)=0$.

Suppose first that $a$ is right-scattered.  The one-step identity for the
Cauchy $\Delta$-integral gives
\[
 0\leq u(\sigma(a))
 \leq \int_a^{\sigma(a)}w(s,u(s))\Deltas
 =\muT(a)w(a,u(a))
 =0.
\]
Thus $u(\sigma(a))=u(a)=0$.  Continuity at a right-scattered point implies
$\Delta$-differentiability, and the difference-quotient formula yields
\[
 u^\Delta(a)=\frac{u(\sigma(a))-u(a)}{\muT(a)}=0.
\]

Suppose instead that $a$ is right-dense.  Fix $\varepsilon>0$.  Uniform
equicontinuity in axiom (ii), applied at $x=0$, and axiom (iii) give an
$\eta>0$ such that
\begin{equation}\label{eq:uniformsmall}
 0\leq y<\eta
 \quad\Longrightarrow\quad
 0\leq w(s,y)<\varepsilon
 \qquad(s\in[a,b]_{\T}).
\end{equation}
Since $u$ is continuous at $a$ and $u(a)=0$, there is $\delta>0$ such that
$u(s)<\eta$ whenever $s\in[a,b]_{\T}$ and $0\leq s-a<\delta$.  Hence, for
$t\in(a,a+\delta)\cap\T$, \eqref{eq:ineq},
\eqref{eq:uniformsmall}, and $\int_a^t1\Deltas=t-a$ imply
\[
 0\leq \frac{u(t)-u(a)}{t-a}
 \leq \frac1{t-a}\int_a^t\varepsilon\Deltas
 =\varepsilon.
\]
Therefore the time-scale limit of the difference quotient at $a$ is zero,
which is precisely $u^\Delta(a)=0$ in the right-dense case.
\end{proof}

\begin{corollary}[Equivalent definition and impossibility result]
Under axioms (i)--(iii), axiom (iv) of Definition~7 is equivalent to the
following apparently stronger statement:
\begin{quote}
$u\equiv0$ is the unique nonnegative continuous function satisfying
\eqref{eq:ineq}, with no endpoint differentiability assumption.
\end{quote}
Consequently, the source paper's proposed example in which a nonzero solution
is excluded only by the condition $u^\Delta(a)=0$ does not exist.
\end{corollary}

\begin{proof}
The unconditional uniqueness statement plainly implies axiom (iv).  The
reverse implication follows from Theorem~\ref{thm:auto}, because every
nonnegative continuous solution of \eqref{eq:ineq} automatically belongs to
the class quantified over in axiom (iv).
\end{proof}

\section*{A Family of Non-Product Examples}

\begin{proposition}[Polynomial Kamke $\Delta$-functions]\label{prop:poly}
Let $m\geq1$, and let
$p_j\in\Rd([a,b]_{\T},[0,\infty))$ for $1\leq j\leq m$.  Then
\begin{equation}\label{eq:poly}
                  w(t,x)=\sum_{j=1}^m p_j(t)x^j
\end{equation}
is a Kamke $\Delta$-function.

If the coefficient vectors
$(p_1(t),\ldots,p_m(t))$ are not all proportional as $t$ varies, then
\eqref{eq:poly} cannot be written as $q(t)h(x)$.
\end{proposition}

\begin{proof}
For fixed $x$, the function $w(\cdot,x)$ is rd-continuous, and
$w(t,0)=0$.  Every $p_j$ is bounded on the compact time-scale interval.
For a fixed $R>0$, the mean-value theorem in the $x$ variable gives, for
$x,y\in[0,R]$,
\[
 |w(t,x)-w(t,y)|
 \leq
 \left(\sum_{j=1}^m j\,\|p_j\|_\infty R^{j-1}\right)|x-y|,
\]
uniformly in $t$.  Thus axioms (i)--(iii) hold.

Let $u$ be a nonnegative continuous function satisfying \eqref{eq:ineq}, and
put $M=\max_{[a,b]_{\T}}u$.  For $0\leq x\leq M$,
\[
 w(t,x)\leq Q_M(t)x,
 \qquad
 Q_M(t):=\sum_{j=1}^m p_j(t)M^{j-1}.
\]
The coefficient $Q_M$ is nonnegative and rd-continuous.  Therefore
\[
 u(t)\leq\int_a^t Q_M(s)u(s)\Deltas.
\]
The time-scale Gronwall inequality (the same result used in Theorem~10 of the
source paper) gives $u(t)\leq0$ on $[a,b]_{\T}$; hence $u\equiv0$.
This proves axiom (iv), indeed in its unconditional form.

Finally, suppose $w(t,x)=q(t)h(x)$.  At any time $t_0$ for which the
coefficient vector is nonzero, $q(t_0)\neq0$, and equality for all $x\geq0$
forces $h$ to be the polynomial $w(t_0,\cdot)/q(t_0)$.  Comparing polynomial
coefficients then shows
\[
 (p_1(t),\ldots,p_m(t))
 =\frac{q(t)}{q(t_0)}(p_1(t_0),\ldots,p_m(t_0))
\]
for every $t$.  Hence all coefficient vectors would be proportional, a
contradiction.
\end{proof}

\begin{corollary}[Explicit non-product example]\label{cor:example}
On every nondegenerate compact time-scale interval $[a,b]_{\T}$,
\[
                         \boxed{\;w(t,x)=x+(t-a)x^2\;}
\]
is a Kamke $\Delta$-function and is not of the form $q(t)h(x)$.
\end{corollary}

\begin{proof}
Apply Proposition~\ref{prop:poly} with $p_1(t)=1$ and $p_2(t)=t-a$.
For a direct factorization check, the $2\times2$ determinant using
$t=a,b$ and $x=1,2$ is
\[
 \det\begin{pmatrix}w(a,1)&w(a,2)\\w(b,1)&w(b,2)\end{pmatrix}
 =\det\begin{pmatrix}1&2\\1+b-a&2+4(b-a)\end{pmatrix}
 =2(b-a)\neq0.
\]
A product $q(t)h(x)$ would give a rank-one matrix and determinant zero.
\end{proof}

\section*{Verification, Scope, and Novelty}

The endpoint proof was audited against the two exhaustive cases in Remark~17
of the source paper.  Its only time-scale identities are the right-scattered
one-step integral formula, the right-dense difference-quotient formula, and
$\int_a^t1\Deltas=t-a$, all stated in the source preliminaries.  The
polynomial example uses the same Gronwall theorem that Oberta invokes for the
linear example in Theorem~10.  No numerical or computer-assisted step enters
the proof.

The result answers precisely the two example directions in the final
Kamke-$\Delta$ paragraph: one direction positively by a large non-product
family, and the other negatively by proving impossibility.  It does not claim
to modify the paper's main existence theorem or resolve its other concluding
remarks.

A bounded novelty check on August 11, 2026 searched the run's registry,
solution, attempt, and proof-gap indexes for arXiv:2512.13602 and the phrases
``Kamke $\Delta$-function,'' ``non-product,'' and ``endpoint
$\Delta$-differentiability.''  Exact-title and close-phrase web searches found
no independent answer.  OpenAlex listed one 2026 citing work, on fractional
retarded dynamic equations, whose indexed title and abstract do not state
either result here.  Novelty confidence is therefore moderate, not
definitive, pending a specialist literature review.

The recommended human-review focus is the quantifier use of uniform
equicontinuity in \eqref{eq:uniformsmall}, the endpoint convention for
$[a,b]_{\T}$, and the application of time-scale Gronwall to $Q_M$.

\begin{thebibliography}{9}
\bibitem{Oberta}
D.~Oberta,
\emph{On the existence of solutions of dynamic equations on time scales in
Banach spaces}, arXiv:2512.13602v2 (2026); published online in
Math. Methods Appl. Sci., DOI: \href{https://doi.org/10.1002/mma.70827}{10.1002/mma.70827}.

\bibitem{BohnerPeterson}
M.~Bohner and A.~Peterson,
\emph{Dynamic Equations on Time Scales: An Introduction with Applications},
Birkh\"auser, Boston, 2001.  See Theorem~6.4 for the Gronwall inequality
used by the source paper and Proposition~\ref{prop:poly}.
\end{thebibliography}

\end{document}
